
\section{Evaluación}

A diferencia de los problemas de Machine Learning clásico, como clasificaciones o regresiones, evaluar los modelos de inteligencia artificial generativa supone un verdadero reto. Eso se debe a que el espacio de \textit{outputs} (respuestas) correctas es infinito. La aproximación en \cite{mi-paper} fue utilizar el método \textit{LLM-as-a-judge}, una solución muy común en el estado del arte dada su simplicidad: consiste en hacer que sea otro modelo de lenguaje el que evalúe las respuestas. Sin embargo, este tipo de evaluaciones son muy poco robustas a alucinaciones técnicas, tienen un sesgo de autorreferencia y tienden a perder consistencia.

Es por ello que \textbf{se ha diseñado una metodología experimental de evaluación}, la cual se detalla en esta sección. Este enfoque, diseñado para optimizar el consumo de recursos operativos, manteniendo la robustez estadística, se compone de dos etapas principales: evaluación del sistema completo y análisis de contribuciones marginales mediante estudio de ablación.

El principal reto de esta metodología experimental, y lo que la diferencia de otras existentes, es que tiene que ser capaz de abordar la evaluación del uso del contexto geoespacial en la elaboración de las respuestas, y no únicamente la calidad de las mismas. No es lo mismo una respuesta que es correcta pero genérica, o incluso aparentemente específica pero generada gracias al conocimiento de base del modelo, que una \textit{personalizada} que se adapte a las características del terreno. Es por ello que se combinan varias etapas, se evalúan diferentes configuraciones, y se proponen métricas específicas. 

\subsection{Selección de parcelas y heterogeneidad}
Se seleccionaron un total de 15 parcelas, categorizadas en cuatro tipos distintos de Uso del Suelo (LU). La distribución de evaluaciones se pondera de acuerdo con la utilidad operativa de la herramienta para cada uso funcional del terreno, según se muestra en \tabref{tab:land_use}. \\

Para garantizar una evaluación libre de sesgos, las parcelas se seleccionaron de forma que representaran un alto grado de heterogeneidad en sus características físicas, incluyendo pendientes variables ($0.5\%$ a $21.6\%$), altitudes, formas irregulares y diversidad en su salud vegetal (rangos de NDVI). Las parcelas se muestran en la \figref{fig:mosaic}.


\begin{figure*}[t!]
    \centering{\includegraphics[width=\textwidth]{figures/mosaic.png}}
    \caption{Selección de parcelas para la evaluación de la herramienta, agrupadas por Land Use. El número de parcelas por cada LU depende de su utilidad funcional en agricultura.}
    \label{fig:mosaic}
\end{figure*}


\begin{table}[H]
    \caption{Distribución de parcelas por tipo de Uso del Suelo (LU)}
    \label{tab:land_use}
    \begin{tabular}{
        >{\raggedright\arraybackslash}p{0.55\columnwidth}
        >{\raggedleft\arraybackslash}p{0.40\columnwidth}
    }
    \toprule
    \textbf{Uso del Suelo (Land Use)} & \textbf{Cantidad de queries} \\
    \midrule
    Árboles Frutales (Fruit Tree - FR) & 7 \\ 
    Pastos (Pasture - PA) & 5 \\ 
    Improductivo (Unproductive - UN) & 2 \\ 
    Forestal (Forest - FO) & 1 \\
    \bottomrule
    \end{tabular}
\end{table}

\subsection{Fase I: Broad Performance Assessment}
Esta fase evalúa la capacidad general del sistema a través de diversos escenarios. Se diseñaron un total de 71 consultas, distribuidas de la siguiente manera:
\begin{itemize}
    \item \textbf{Parcelas Productivas (FR y PA):} 5 consultas por parcela, centradas en riego, fertilización, gestión de plagas y el ROI de la inversión.
    \item \textbf{Parcelas No Productivas (UN y FO):} 4 consultas por parcela, que consisten en 2 consultas funcionales y 2 consultas de pruebas negativas (\textit{Negative Testing Queries}, escenarios donde el sistema debe identificar la no aptitud).
\end{itemize}

\subsection{Fase II: Ablation Study (Validación de Arquitectura)}
Para cuantificar el impacto y aportación marginal de cada módulo de la arquitectura, se realizó un estudio de ablación utilizando un \textbf{diseño factorial}. Los estudios de ablación consisten en el análisis individual de los distintos subconjuntos (configuraciones) de un sistema, obtenidos mediante la activación y desactivación sistemática de sus componentes.  
En este caso, al contar con tres módulos (Ortofoto, Satélite y RAG, mostrado anteriormente en \figref{fig:arch}), el diseño factorial de experimentos da como resultado $2^3 = 8$ configuraciones posibles.

\begin{itemize}
    \item \textbf{Configuraciones:} 8 combinaciones modulares + 1 línea base Zero-Shot (el LLM sin contexto de la parcela) = 9 configuraciones en total, recogidas en \tabref{tab:ablation_configs}.
    \item \textbf{Subconjunto de muestra:} 1 parcela representativa por cada categoría de Uso del Suelo.
    \item \textbf{Volumen de consultas:} 2 consultas específicas por parcela a lo largo de las 9 configuraciones $= \mathbf{72 \text{ consultas}}$.
\end{itemize}
\begin{table}[H]
    \caption{Configuraciones modulares del estudio de ablación.}
    \label{tab:ablation_configs}
    \begin{tabular}{
        >{\raggedright\arraybackslash}p{0.28\columnwidth}
        >{\centering\arraybackslash}p{0.07\columnwidth}
        >{\centering\arraybackslash}p{0.07\columnwidth}
        >{\centering\arraybackslash}p{0.07\columnwidth}
        >{\centering\arraybackslash}p{0.08\columnwidth}
        >{\centering\arraybackslash}p{0.17\columnwidth}
        >{\raggedleft\arraybackslash}p{0.09\columnwidth}
    }
    \toprule
    \textbf{Módulos} & \textbf{R} & \textbf{O} & \textbf{S} & \textbf{PD} & \textbf{Tool Activa} & \textbf{Código} \\
    \midrule
    Baseline & 0 & 0 & 0 & 1 & 1 & 0 \\
    B + SAT & 0 & 0 & 1 & 1 & 1 & 1 \\
    B + ORTHO & 0 & 1 & 0 & 1 & 1 & 2 \\
    B + ORTHO + SAT & 0 & 1 & 1 & 1 & 1 & 3 \\
    B + RAG & 1 & 0 & 0 & 1 & 1 & 4 \\
    B + RAG + SAT & 1 & 0 & 1 & 1 & 1 & 5 \\
    B + RAG + ORTHO & 1 & 1 & 0 & 1 & 1 & 6 \\
    Complete & 1 & 1 & 1 & 1 & 1 & 7 \\
    ZERO SHOT & 0 & 0 & 0 & 0 & 1 & ZS \\
    \bottomrule
    \end{tabular}
    \tabletext{Nota: PD = Módulo SIGPAC. El valor Tool activa siempre está activo, ya que es la encargada de identificar los códigos de parcela como tales. }
\end{table}

\subsection{Resumen del volumen de evaluación}
Los resultados combinados proporcionan un conjunto de 143 evaluaciones totales (\tabref{tab:evaluation_summary}). En el Anexo A, se incluye el enlace al conjunto de datos elaborado para esta evaluación, donde se pueden consultar todas las queries, respuestas y contexto. 

\begin{table}[H]
    \caption{Resumen del volumen total de consultas en la evaluación.}
    \label{tab:evaluation_summary}
    \begin{tabular}{
        >{\raggedright\arraybackslash}p{0.35\columnwidth}
        >{\raggedright\arraybackslash}p{0.22\columnwidth}
        >{\raggedleft\arraybackslash}p{0.35\columnwidth}
    }
    \toprule
    \textbf{Fase de evaluación} & \textbf{Metodología} & \textbf{Num. Consultas} \\
    \midrule
    Broad Assessment & Secuencial & 71 \\
    Ablation Study & Simultáneo & 72 \\
    \midrule
    \textbf{Total} & & \textbf{143} \\
    \bottomrule
    \end{tabular}
\end{table}

\subsection{Sujetos de evaluación y mitigación de sesgos}
Para garantizar una validación equilibrada, mitigar sesgos individuales y capturar un rango de perspectivas amplio, el proceso de evaluación es llevado a cabo de forma independiente por \textbf{cuatro evaluadores} diferenciados:

\begin{itemize}
    \item \textbf{Evaluadores Humanos (2):}
    \begin{itemize}
        \item \textit{Evaluador Experto:} Doctor en Biología.
        \item \textit{Evaluador de Control:} Perfil no especialista.
    \end{itemize}
    \item \textbf{LLM-as-a-judge (2):}
    \begin{itemize}
        \item \textit{Modelo Propietario (Gemini)}: Gemini3.1 Flash Lite 
        \item \textit{Modelo Open Source (Qwen)}: Qwen3.5-397B-A17B-Instruct
    \end{itemize}
\end{itemize}

\begin{figure*}[t!]
    \centering{\includegraphics[width=\textwidth]{figures/configurations.png}}
    \caption{Resultados preliminares de la evaluación en la fase de \textit{Ablation Study} para los 4 evaluadores y las 5 métricas agrupados por configuración. Se pueden apreciar diferencias según los módulos activos. }
    \label{fig:results}
\end{figure*}

\subsection{Marco de métricas multidimensionales}

La calidad de las respuestas se mide según un baremo multidimensional formado por cinco métricas, divididas en las que son puntuadas teniendo en cuenta el contexto y las que no, tal y como muestra la \tabref{tab:evaluation_metrics}. Las puntuaciones se llevan a cabo en la escala \textit{Likert} (valores enteros entre 1 y 5). 

\begin{table}[H]
    \caption{Marco de métricas de evaluación multidimensional.}
    \label{tab:evaluation_metrics}
    \begin{tabular}{
        >{\raggedright\arraybackslash}p{0.45\columnwidth}
        >{\raggedleft\arraybackslash}p{0.50\columnwidth}
    }
    \toprule
    \textbf{Métrica} & \textbf{Tipo} \\
    \midrule
    Correctness & Without Context \\
    Completeness & Without Context \\
    Relevance & Without Context \\
    Utility & Without Context \\
    \midrule
    Grounding & With Context \\
    \bottomrule
    \end{tabular}
    \tabletext{Nota: La métrica \textit{Grounding} mdie el grado en el que las respuestas están derivadas del contexto obtenido, y no son fruto de conocimiento genérico, por lo que es la única que necesita el contexto para ser evaluada. La métrica \textit{Correctness} no aplica para el evaluador humano de control.}
\end{table}

\subsection{Resultados preliminares y discusión}

En esta sección se presentan algunos resultados preliminares de la evaluación. Cabe mencionar que las gráficas incluidas aquí son únicamente una primera visión que arroja luz sobre la evaluación del sistema; los análisis estadísticos más avanzados se reservan para una futura publicación científica. En primer lugar, la \figref{fig:results} contiene las puntuaciones promedio de cada uno de los 4 evaluadores en las 5 métricas, agrupadas por configuración. De esta gráfica pueden sacarse multitud de observaciones, siendo las más destacadas las relativas a la correlación entre evaluadores; y especialmente en referencia a las aportaciones marginales de los módulos. Recordando la \tabref{tab:ablation_configs}, vemos que las configuraciones 2 y 3, además de la 6 y 7, son las que incluyen el módulo de ortofoto; y en todas las métricas se observa una mayor puntuación promedio en dichos casos, por lo que se puede deducir que la aportación marginal del módulo ortofoto es significativa. 

\begin{figure}[H]
    \centering
    \includegraphics[width=\columnwidth]{figures/broad_bars.png }
    \caption{Valores promedio para los 4 evaluadores y las 5 métricas en la fase \textit{Broad} (sistema completo).}
    \label{fig:broad}
\end{figure}

De la misma manera, la tendencia en todas las métricas es ascendente con el eje x (los valores son más bajos en las configuraciones de la izquierda), viendo como, por ejemplo, la utilidad del sistema en la configuración \textit{Zero Shot}, es inferior a 1/5 para la mayoría de evaluadores. Dado que esta configuración corresponde a un modelo LLM estándar sin ampliación de contexto, y estando ordenadas las configuraciones progresivamente con los módulos activos, se puede concluir que \textbf{las respuestas, tal y como se esperaba, son mejores a medida que se añade información geoespacial del terreno.} Este resultado, por consiguiente, sirve para verificar la hipótesis inicial y validar la arquitectura planteada. 

De manera más genérica, la \figref{fig:broad} representa los valores promedios de las métricas sobre las evaluaciones de la etapa \textit{Broad}, es decir, sobre el sistema completo. Para todos los evaluadores, las métricas resultantes promedian valores muy cercanos a 5/5, lo que implica que el sistema ha respondido con alto grado de éxito las consultas evaluadas. Estos resultados preliminares sirven como evidencia de que el sistema propuesto es capaz de generar respuestas de calidad, completas, útiles, y en general exitosas en el sentido amplio. 

% La metodología experimental diseñada, que puede ser replicada en otros estudios, evalúa tanto la totalidad del sistema como sus partes por separado, brindando la capacidad de medir las aportaciones de cada módulo.


\section{Conclusiones}

El desarrollo de este proyecto se proponía como objetivo \textbf{crear un puente entre la información geoespacial}, compleja y no estructurada, \textbf{y el usuario final.} A través de modelos de lenguaje multimodaldes de código abierto, y una arquitectura modular que integra diversas fuentes de datos multidominio y las transforma a lenguaje natural, se propone una herramienta que no solo da utilidad a recursos públicos de gran potencial, sino que además democratiza el aprovechamiento de la inteligencia geoespacial, eliminando la barrera técnica de entrada. La utilidad de esta herramienta comprende desde la evaluación de las características de una parcela como la irrigación o la calidad del suelo, hasta la viabilidad y comparación de distintas plantaciones, pasando por cualquier consulta relacionada con agriculura, gestión del terreno o evaluaciones de suelos. Esto permite al usuario tener un asistente virtual experto que conoce su terreno, así como el contexto geoespacial, técnico, legal, etc. Esto supondría un gran impacto positivo en la sostenibilidad, eficiencia y competitividad de las explotaciones agrarias de regiones rurales como Asturias. El sistema ha sido validado mediante una metodología de evaluación multi-etapa que incluye análisis de aportaciones marginales de los módulos mediante un estudio de ablación. Los resultados obtenidos demuestran la capacidad del sistema para generar respuestas de calidad, así como la mejoría que representa este contexto geoespacial añadido. Tanto la herramienta validada como la metodología de evaluación propuesta y el conjunto de datos construido para tal fin, suponen contribuciones adicionales a la comunidad científica. 
